\chapter{2007年湖北卷（文）}

\section{单选题}

\begin{problem}
\(\tan 690^{\circ}\) 的值为（\quad）

\choices
    {\( -\frac{\sqrt{3}}{3} \)}
    {\( \frac{\sqrt{3}}{3} \)}
    {\(\sqrt{3}\)}
    {\(-\sqrt{3}\)}
\end{problem}

\begin{answer}
A.
\end{answer}

\begin{solution}
利用正切函数的周期性，
\[
690^{\circ}=720^{\circ}-30^{\circ}
\]
所以
\[
\tan 690^{\circ}=\tan(-30^{\circ})=-\tan 30^{\circ}=-\frac{\sqrt{3}}{3}
\]
故选项 A 正确.
\end{solution}

\begin{problem}
如果 \( U = \{x \mid x \text{ 是小于 } 9 \text{ 的正整数}\} \)，\( A = \{1, 2, 3, 4\} \)，\( B = \{3, 4, 5, 6\} \)，那么 \( \complement_U A \cap \complement_U B = \)（\quad）

\choices
    {\( \{1, 2\} \)}
    {\( \{3, 4\} \)}
    {\( \{5, 6\} \)}
    {\(\{7,8\}\)}
\end{problem}

\begin{answer}
D.
\end{answer}

\begin{solution}
全集为
\[
U=\{1,2,3,4,5,6,7,8\}
\]
因此
\(\complement_U A=\{5,6,7,8\}\)，\(\complement_U B=\{1,2,7,8\}\).
两者的交集为
\[
\complement_U A\cap\complement_U B=\{7,8\}
\]
故选项 D 正确.
\end{solution}

\begin{problem}
如果 \(\left(3x^{2} - \frac{2}{x^{3}}\right)^{n}\) 的展开式中含有非零常数项，则正整数 \(n\) 的最小值为（\quad）

\choices
    {\(10\)}
    {\(6\)}
    {\(5\)}
    {\(3\)}
\end{problem}

\begin{answer}
C.
\end{answer}

\begin{solution}
在展开式中，取含有 \(k\) 个因式 \(-2x^{-3}\) 的项，所得项为
\[
\mathrm C_n^k(3x^2)^{n-k}\left(-\frac{2}{x^3}\right)^k
\]
其幂的指数为
\[
2(n-k)-3k=2n-5k
\]
含有非零常数项的条件是 \(2n-5k=0\). 因为 \(2\) 与 \(5\) 互质，所以 \(k\) 必须是偶数，\(n\) 必须是 \(5\) 的倍数. 取最小的正整数 \(k=2\)，得到 \(n=5\)，且此时 \(k\leqslant n\) 满足选取条件. 故选项 C 正确.
\end{solution}

\begin{problem}
函数 \( y = \frac{2^x + 1}{2^x - 1} \)（\(x < 0\)）的反函数是（\quad）

\choices
    {\( y = \log_2\frac{x + 1}{x - 1} \)（\(x < - 1\)）}
    {\( y = \log_2\frac{x + 1}{x - 1} \)（\(x > 1\)）}
    {\( y = \log_2\frac{x - 1}{x + 1} \)（\(x < - 1\)）}
    {\( y = \log_2\frac{x - 1}{x + 1} \)（\(x > 1\)）}
\end{problem}

\begin{answer}
A.
\end{answer}

\begin{solution}
令
\(y=\frac{2^x+1}{2^x-1}\)，\(x<0\).
移项得
\[
(y-1)2^x=y+1
\]
从而
\(2^x=\frac{y+1}{y-1}\)，\(x=\log_2\frac{y+1}{y-1}\).
当 \(x<0\) 时，\(0<2^x<1\)，所以
\[
0<\frac{y+1}{y-1}<1
\]
解得 \(y<-1\). 将反函数的自变量记为 \(x\)，因变量记为 \(y\)，得
\(y=\log_2\frac{x+1}{x-1}\)，\(x<-1\).
故选项 A 正确.
\end{solution}

\begin{problem}
在棱长为 \(1\) 的正方体 \(ABCD - A_{1}B_{1}C_{1}D_{1}\) 中，\(E\)，\(F\) 分别为棱 \(AA_{1}\)，\(BB_{1}\) 的中点，\(G\) 为棱 \(A_{1}B_{1}\) 上的一点，且 \(A_{1}G = \lambda\)（\(0 \leqslant \lambda \leqslant 1\)），则点 \(G\) 到平面 \(D_{1}EF\) 的距离为（\quad）

\bitmapfigure[max width=0.34\linewidth]{img/2007/hubei_liberal/q5_fig1.png}

\choices
    {\(\sqrt{3}\)}
    {\(\frac{\sqrt{2}}{2}\)}
    {\(\frac{\sqrt{2}\lambda}{3}\)}
    {\(\frac{\sqrt{5}}{5}\)}
\end{problem}

\begin{answer}
D.
\end{answer}

\begin{solution}
建立直角坐标系，取
\(A(0,0,0)\)，\(\ B(1,0,0)\)，\(\ D(0,1,0)\)，\(\ A_1(0,0,1)\).
则
\(D_1(0,1,1)\)，\(E\left(0,0,\frac12\right)\)，\(F\left(1,0,\frac12\right)\)，\(G(\lambda,0,1)\).
平面 \(D_1EF\) 内有两个相交向量
\(\overrightarrow{D_1E}=\left(0,-1,-\frac12\right)\)，\(\overrightarrow{D_1F}=\left(1,-1,-\frac12\right)\).
取其法向量的一个倍数为 \((0,-1,2)\)，故平面方程为
\[
-y+2z-1=0
\]
代入 \(G\) 的坐标，得点 \(G\) 到该平面的距离
\[
d=\frac{|0+2-1|}{\sqrt{0^2+(-1)^2+2^2}}=\frac{1}{\sqrt5}=\frac{\sqrt5}{5}
\]
这个结果与 \(\lambda\) 无关，故选项 D 正确.
\end{solution}

\begin{problem}
为了了解某学校学生的身体发育情况，抽查了该校 \(100\) 名高中男生的体重情况，根据所得数据画出样本的频率分布直方图如图所示，根据此图，估计该校 \(2000\) 名高中男生中体重大于 \(70.5\) 公斤的人数为（\quad）

\bitmapfigure[max width=0.95\linewidth]{img/2007/hubei_liberal/q6_fig1.png}

\choices
    {\(300\)}
    {\(360\)}
    {\(420\)}
    {\(450\)}
\end{problem}

\begin{answer}
B.
\end{answer}

\begin{solution}
由频率分布直方图可读出，体重在 \(70.5\) 至 \(72.5\)、\(72.5\) 至 \(74.5\)、\(74.5\) 至 \(76.5\) 公斤三个组的频率密度分别为 \(0.04\)、\(0.035\)、\(0.015\)，每组组距均为 \(2\). 因而体重大于 \(70.5\) 公斤的样本频率为
\[
2(0.04+0.035+0.015)=0.18
\]
用样本频率估计总体频率，\(2000\) 名男生中体重大于 \(70.5\) 公斤的人数约为
\[
2000\times0.18=360
\]
故选项 B 正确.
\end{solution}

\begin{problem}
将 \(5\) 本不同的书全发给 \(4\) 名同学，每名同学至少有一本书的概率是（\quad）

\choices
    {\( \frac{15}{64} \)}
    {\( \frac{15}{128} \)}
    {\( \frac{24}{125} \)}
    {\( \frac{48}{125} \)}
\end{problem}

\begin{answer}
A.
\end{answer}

\begin{solution}
把每本不同的书发给一名同学，每本书都有 \(4\) 种去向，所以全部发法共有 \(4^5\) 种.

用容斥原理计算每名同学至少得到一本书的发法数：
\[
4^5-\mathrm C_4^1 3^5+\mathrm C_4^2 2^5-\mathrm C_4^3 1^5
=1024-972+192-4=240
\]
因此所求概率为
\[
\frac{240}{4^5}=\frac{240}{1024}=\frac{15}{64}
\]
故选项 A 正确.
\end{solution}

\begin{problem}
由直线 \( y = x + 1 \) 上的一点向圆 \( (x - 3)^2 + y^2 = 1 \) 引切线，则切线长的最小值为（\quad）

\choices
    {\(1\)}
    {\(2 \sqrt{2}\)}
    {\(\sqrt{7}\)}
    {\(3\)}
\end{problem}

\begin{answer}
C.
\end{answer}

\begin{solution}
圆心为 \(O(3,0)\)，半径为 \(1\). 设直线 \(y=x+1\) 上的动点为 \(P\)，从 \(P\) 引圆的切线，切线长 \(L\) 满足
\[
L^2=PO^2-1
\]
因此要使 \(L\) 最小，只需使 \(PO\) 最小，即取圆心到直线的距离. 直线化为 \(x-y+1=0\)，故
\[
PO_{\min}=\frac{|3-0+1|}{\sqrt{1^2+(-1)^2}}=2\sqrt2
\]
于是
\[
L_{\min}=\sqrt{(2\sqrt2)^2-1}=\sqrt7
\]
故选项 C 正确.
\end{solution}

\begin{problem}
设 \(\bs{a}=(4,3)\)，\(\bs{a}\) 在 \(\bs{b}\) 上的投影为 \( \frac{5\sqrt{2}}{2} \)，\(\bs{b}\) 在 \(x\) 轴上的投影为 \(2\)，且 \(\left|\bs{b}\right|\leqslant14\)，则 \(\bs{b}\) 为（\quad）

\choices
    {\((2,14)\)}
    {\( \left(2,-\frac{2}{7}\right) \)}
    {\( \left(-2,\frac{2}{7}\right) \)}
    {\((2,8)\)}
\end{problem}

\begin{answer}
B.
\end{answer}

\begin{solution}
由 \(\bs{b}\) 在 \(x\) 轴上的投影为 \(2\)，设
\[
\bs b=(2,t)
\]
向量 \(\bs{a}\) 在 \(\bs{b}\) 上的投影为
\[
\frac{\bs a\cdot\bs b}{|\bs b|}=\frac{8+3t}{\sqrt{4+t^2}}=\frac{5\sqrt2}{2}
\]
平方并整理，得
\[
(8+3t)^2=\frac{25}{2}(4+t^2)
\]
即
\[
7t^2-96t-28=0
\]
所以 \(t=14\) 或 \(t=-2/7\). 当 \(t=14\) 时，
\[
|\bs b|=\sqrt{2^2+14^2}>14
\]
不符合条件；当 \(t=-2/7\) 时原投影等式成立. 因而
\[
\bs b=\left(2,-\frac27\right)
\]
故选项 B 正确.
\end{solution}

\begin{problem}
已知 \(p\) 是 \(r\) 的充分条件而不是必要条件，\(q\) 是 \(r\) 的充分条件，\(s\) 是 \(r\) 的必要条件，\(q\) 是 \(s\) 的必要条件. 现有下列命题：

\begin{circlelist}
\item \(s\) 是 \(q\) 的充要条件；
\item \(p\) 是 \(q\) 的充分条件而不是必要条件；
\item \(r\) 是 \(q\) 的必要条件而不是充分条件；
\item \(\neg p\) 是 \(\neg s\) 的必要条件而不是充分条件；
\item \(r\) 是 \(s\) 的充分条件而不是必要条件.
\end{circlelist}

则正确命题的序号是（\quad）

\choices
    {\(\circled{1}\circled{4}\circled{5}\)}
    {\(\circled{1}\circled{2}\circled{4}\)}
    {\(\circled{2}\circled{3}\circled{5}\)}
    {\(\circled{2}\circled{4}\circled{5}\)}
\end{problem}

\begin{answer}
B.
\end{answer}

\begin{solution}
题设条件可写成逻辑关系
\(p\Rightarrow r\)，\(r\Rightarrow s\)，\(s\Rightarrow q\).
又因为 \(q\) 是 \(r\) 的充分条件，所以 \(q\Rightarrow r\). 因而
\(r\Longleftrightarrow s\Longleftrightarrow q\)，\(p\Rightarrow q\).

命题 1 正确，因为 \(s\) 与 \(q\) 等价；命题 2 正确，因为 \(p\Rightarrow q\)，且题设已知存在 \(r\land\neg p\)，由 \(r\Leftrightarrow q\) 可知 \(p\) 不是 \(q\) 的必要条件；命题 3 错误，因为 \(r\) 不仅是 \(q\) 的必要条件，也是其充分条件；命题 4 正确，因为 \(p\Rightarrow s\) 的逆否命题是 \(\neg s\Rightarrow\neg p\)，且由 \(r\land\neg p\) 及 \(r\Leftrightarrow s\) 可知存在 \(s\land\neg p\)，所以 \(\neg p\) 不是 \(\neg s\) 的充分条件；命题 5 错误，因为 \(r\) 与 \(s\) 等价，不是“充分条件而不是必要条件”.

所以正确命题的序号为 \(1\)，\(2\)，\(4\)，故选项 B 正确.
\end{solution}

\section{填空题}

\begin{problem}
设变量 \(x\)，\(y\) 满足约束条件 \( \begin{cases}
x-y+3&\geqslant0,\\
x+y&\geqslant0,\\
-2&\leqslant x\leqslant3,
\end{cases} \)，则目标函数 \(2x+y\) 的最小值为\fillinblank{}.
\end{problem}

\begin{answer}
\(-\frac{3}{2}\).
\end{answer}

\begin{solution}
约束条件等价于
\(y\leqslant x+3\)，\(y\geqslant-x\)，\(-2\leqslant x\leqslant3\).
要存在满足前两个不等式的 \(y\)，必须有
\(-x\leqslant x+3\)，\(x\geqslant-\frac32\).
对固定的 \(x\)，目标函数 \(2x+y\) 随 \(y\) 增大而增大，所以取允许的最小值 \(y=-x\). 此时
\[
2x+y=x
\]
可行的 \(x\) 范围为 \([-\frac32,3]\)，故目标函数的最小值为
\[
-\frac32
\]
\end{solution}

\begin{problem}
过双曲线 \( \frac{x^{2}}{4} - \frac{y^{2}}{3} = 1 \) 左焦点 \(F_{1}\) 的直线交双曲线的左支于 \(M\)，\(N\) 两点，\(F_{2}\) 为其右焦点，则 \( \left|MF_{2}\right| + \left|NF_{2}\right| - \left|MN\right| \) 的值为\fillinblank{}.
\end{problem}

\begin{answer}
\(8\).
\end{answer}

\begin{solution}
双曲线的实半轴为 \(a=2\)，焦半距为
\[
c=\sqrt{a^2+b^2}=\sqrt{4+3}=\sqrt7
\]
所以左、右焦点分别为 \(F_1(-\sqrt7,0)\)、\(F_2(\sqrt7,0)\). 对左支上任一点 \(X\)，双曲线的定义给出
\[
|XF_2|-|XF_1|=2a=4
\]
下面说明 \(F_1\) 位于 \(M\)，\(N\) 之间. 若直线不是竖直直线，设其方程为 \(y=k(x+\sqrt7)\)，并令 \(u=x+\sqrt7\). 代入双曲线方程得
\[
(3-4k^2)u^2-6\sqrt7u+9=0
\]
两个交点都在左支上，故 \(u\leqslant\sqrt7-2\). 若 \(3-4k^2>0\)，则两根为正且根的和
\[
\frac{6\sqrt7}{3-4k^2}\geqslant2\sqrt7>2(\sqrt7-2)
\]
不可能两根都不超过 \(\sqrt7-2\)；若 \(3-4k^2=0\)，方程只有一个有限根. 因而 \(3-4k^2<0\)，两根之积为
\[
\frac{9}{3-4k^2}<0
\]
所以 \(u=0\)（即 \(F_1\)）位于两交点之间. 竖直直线 \(x=-\sqrt7\) 与双曲线交于 \(y=\pm\frac32\)，结论同样成立. 因此
\[
|MN|=|MF_1|+|NF_1|
\]
于是
\[
|MF_2|+|NF_2|-|MN|
=(|MF_1|+4)+(|NF_1|+4)-(|MF_1|+|NF_1|)=8
\]
\end{solution}

\begin{problem}
已知函数 \( y = f(x) \) 的图象在点 \( M(1, f(1)) \) 处的切线方程是 \( y = \frac{1}{2} x + 2 \)，则 \( f(1) + f'(1) = \)\fillinblank{}.
\end{problem}

\begin{answer}
\(3\).
\end{answer}

\begin{solution}
切线在 \(M(1,f(1))\) 处，所以点 \(M\) 在直线 \(y=\frac12x+2\) 上，得到
\[
f(1)=\frac12+2=\frac52
\]
切线斜率等于函数在该点的导数，故
\[
f'(1)=\frac12
\]
因此
\[
f(1)+f'(1)=\frac52+\frac12=3
\]
\end{solution}

\begin{problem}
某男运动员在三分线投球的命中率是 \(\frac{1}{2}\)，他投球 \(10\) 次，恰好投进 \(3\) 个球的概率是\fillinblank{}. （用数值作答）
\end{problem}

\begin{answer}
\(\frac{15}{128}\).
\end{answer}

\begin{solution}
设投进球的次数为 \(X\)，则 \(X\sim B\left(10,\frac12\right)\). 恰好投进 \(3\) 个球的概率为
\[
P(X=3)=\mathrm C_{10}^3\left(\frac12\right)^3\left(1-\frac12\right)^7
=120\left(\frac12\right)^{10}
=\frac{120}{1024}=\frac{15}{128}
\]
\end{solution}

\begin{problem}
为了预防流感，某学校对教室用药熏消毒法进行消毒. 已知药物释放过程中，室内每立方米空气中的含药量 \(y\)（毫克）与时间 \(t\)（小时）成正比；药物释放完毕后，\(y\) 与 \(t\) 的函数关系式为 \(y = \left(\frac{1}{16}\right)^{t - a}\)（\(a\) 为常数），如图所示.

\begin{enumerate}
\item 从药物释放开始，每立方米空气中的含药量 \(y\)（毫克）与时间 \(t\)（小时）之间的函数关系式为\fillinblank{}；
\item 据测定，当空气中每立方米的含药量降低到 \(0.25\) 毫克以下时，学生方可进教室，那么药物释放开始，至少需要经过\fillinblank{}小时后，学生才能回到教室.
\end{enumerate}

\bitmapfigure[max width=0.34\linewidth]{img/2007/hubei_liberal/q15_fig1.png}
\end{problem}

\begin{answer}
\(
y=\begin{cases}
10t,&0\leqslant t\leqslant\frac{1}{10},\\
\displaystyle\left(\frac{1}{16}\right)^{t-\frac{1}{10}},&t>\frac{1}{10},
\end{cases}
\)，\(0.6\).
\end{answer}

\begin{solution}
\begin{enumerate}
\item
药物释放过程中 \(y\) 与 \(t\) 成正比，设 \(y=kt\). 由图知 \(t=0.1\) 时 \(y=1\)，所以 \(k=10\)，释放阶段为
\(y=10t\)，\(0\leqslant t\leqslant\frac1{10}\).
释放完毕后函数为
\[
y=\left(\frac1{16}\right)^{t-a}
\]
图像在 \(t=0.1\) 处连续且函数值为 \(1\)，故
\[
1=\left(\frac1{16}\right)^{\frac1{10}-a}
\]
由于底数不为 \(1\)，所以 \(a=1/10\). 因而完整函数为
\[
y=\begin{cases}
10t,&0\leqslant t\leqslant\frac1{10},\\
\displaystyle\left(\frac1{16}\right)^{t-\frac1{10}},&t>\frac1{10}.
\end{cases}
\]

\item
当 \(y\leqslant0.25=1/4\) 时，必在释放完毕后. 由
\[
\left(\frac1{16}\right)^{t-\frac1{10}}\leqslant\frac14
=\left(\frac1{16}\right)^{1/2}
\]
且 \(0<1/16<1\)，可得
\(t-\frac1{10}\geqslant\frac12\)，\(t\geqslant\frac35=0.6\).
所以至少经过 \(0.6\) 小时.
\end{enumerate}
\end{solution}

\section{解答题}

\begin{problem}
已知函数 \(f(x) = 2\sin^2\left(\frac{\pi}{4} + x\right) - \sqrt{3}\cos 2x\)，\(x \in \left[\frac{\pi}{4}, \frac{\pi}{2}\right]\).
\begin{enumerate}
\item 求 \(f(x)\) 的最大值和最小值.
\item 若不等式 \( \left|f(x)-m\right|<2 \) 在 \( x\in\left[\frac{\pi}{4},\frac{\pi}{2}\right] \) 上恒成立，求实数 \(m\) 的取值范围.
\end{enumerate}
\end{problem}

\begin{solution}
\begin{enumerate}
\item
由
\[
2\sin^2\left(\frac{\pi}{4}+x\right)
=1-\cos\left(\frac{\pi}{2}+2x\right)
=1+\sin2x
\]
得
\[
f(x)=1+\sin2x-\sqrt3\cos2x
=1+2\sin\left(2x-\frac{\pi}{3}\right)
\]
当 \(x\in[\frac{\pi}{4},\frac{\pi}{2}]\) 时，
\[
2x-\frac{\pi}{3}\in\left[\frac{\pi}{6},\frac{2\pi}{3}\right]
\]
在此区间内，正弦函数的最大值为 \(1\)，最小值为 \(\frac12\)，故
\(f_{\max}=1+2=3\)，\(f_{\min}=1+2\cdot\frac12=2\).
\item
因此 \(|f(x)-m|<2\) 对区间内一切 \(x\) 恒成立，当且仅当
\(|2-m|<2\)，\(|3-m|<2\).
这两个不等式分别给出 \(0<m<4\) 和 \(1<m<5\)，所以
\[
1<m<4
\]
\end{enumerate}
\end{solution}

\begin{problem}
如图，在三棱锥 \(V - ABC\) 中，\(VC \perp\) 底面 \(ABC\)，\(AC \perp BC\)，\(D\) 是 \(AB\) 的中点，且 \(AC = BC = a\)，\(\angle VDC = \theta\)（\(0 < \theta < \frac{\pi}{2}\)）.

\begin{enumerate}
\item 求证：平面 \(VAB\) \(\perp\) 平面 \(VCD\).
\item 试确定角 \(\theta\) 的值，使得直线 \(BC\) 与平面 \(VAB\) 所成的角为 \(\frac{\pi}{6}\).
\end{enumerate}

\bitmapfigure[max width=0.34\linewidth]{img/2007/hubei_liberal/q17_fig1.png}
\end{problem}

\begin{solution}
\begin{enumerate}
\item
取直角坐标系，使
\(C(0,0,0)\)，\(A(a,0,0)\)，\(B(0,a,0)\)，\(V(0,0,h)\)
其中 \(h>0\). 则
\[
D\left(\frac a2,\frac a2,0\right)
\]
因为 \(AC=BC\)，三角形 \(ABC\) 是以 \(C\) 为直角顶点的等腰直角三角形. \(CD\) 是底边 \(AB\) 的中线，所以同时也是高，从而
\[
CD\perp AB
\]
又因为 \(VC\perp\) 底面 \(ABC\)，所以 \(VC\perp AB\). 直线 \(CD\)，\(VC\) 相交且都在平面 \(VCD\) 内，故
\[
AB\perp\text{平面 }VCD
\]
而 \(AB\subset\text{平面 }VAB\)，所以
\[
\text{平面 }VAB\perp\text{平面 }VCD
\]
\item
在直角三角形 \(VCD\) 中，\(\angle VCD=\frac{\pi}{2}\)，故
\(CD=\frac{AB}{2}=\frac{a}{\sqrt2}\)，\(h=CD\tan\theta=\frac{a}{\sqrt2}\tan\theta\).
平面 \(VAB\) 的一个法向量为
\[
\overrightarrow{AB}\times\overrightarrow{AV}
=(-a,a,0)\times(-a,0,h)=(ah,ah,a^2)
\]
可取其方向向量为 \((h,h,a)\). 直线 \(BC\) 的方向向量为 \((0,-a,0)\). 设直线 \(BC\) 与平面 \(VAB\) 所成的角为 \(\varphi\)，则
\[
\sin\varphi
=\frac{|(0,-a,0)\cdot(h,h,a)|}
{a\sqrt{2h^2+a^2}}
=\frac{h}{\sqrt{2h^2+a^2}}
\]
令 \(\varphi=\frac{\pi}{6}\)，得到
\(\frac{h}{\sqrt{2h^2+a^2}}=\frac12\)，\(\Longrightarrow\)，\(4h^2=2h^2+a^2\)，\(\Longrightarrow\)，\(h=\frac{a}{\sqrt2}\).
与 \(h=\frac{a}{\sqrt2}\tan\theta\) 联立，得 \(\tan\theta=1\). 由于 \(0<\theta<\frac{\pi}{2}\)，所以
\[
\theta=\frac{\pi}{4}
\]
\end{enumerate}
\end{solution}

\begin{problem}
某商品每件成本 \(9\) 元. 售价 \(30\) 元，每星期卖出 \(432\) 件. 如果降低价格，销售量可以增加，且每星期多卖出的商品件数与商品单价的降低值 \(x\)（单位：元，\(0 \leqslant x \leqslant 30\)）的平方成正比. 已知商品单价降低 \(2\) 元时，一星期多卖出 \(24\) 件.
\begin{enumerate}
\item 将一个星期的商品销售利润表示成 \(x\) 的函数.
\item 如何定价才能使一个星期的商品销售利润最大？
\end{enumerate}
\end{problem}

\begin{solution}
\begin{enumerate}
\item
设每星期多卖出 \(kx^2\) 件. 由 \(x=2\) 时多卖出 \(24\) 件，得
\(4k=24\)，\(k=6\).
降价 \(x\) 元时，售价为 \(30-x\)，每件利润为 \(21-x\)，销售量为 \(432+6x^2\). 因而一个星期的利润为
\(P(x)=(21-x)(432+6x^2)\)，\(0\leqslant x\leqslant30\).
\item
展开得
\[
P(x)=-6x^3+126x^2-432x+9072
\]
所以
\[
P'(x)=-18x^2+252x-432=-18(x-2)(x-12)
\]
因此 \(P\) 在 \([0,2]\) 上递减，在 \([2,12]\) 上递增，在 \([12,30]\) 上递减. 比较端点和驻点（其中 \(P(2)=8664<P(0)\)）：
\(P(0)=9072\)，\(P(12)=11664\)，\(P(30)=-52488\).
故最大利润在 \(x=12\) 时取得，即售价应定为
\[
30-12=18\text{ 元}
\]
\end{enumerate}
\end{solution}

\begin{problem}
设二次函数 \( f(x) = x^{2} + ax + a \)，方程 \( f(x) - x = 0 \) 的两根 \( x_{1} \) 和 \( x_{2} \) 满足 \( 0 < x_{1} < x_{2} < 1 \).
\begin{enumerate}
\item 求实数 \(a\) 的取值范围.
\item 试比较 \( f(0)f(1) - f(0) \) 与 \( \frac{1}{16} \) 的大小，并说明理由.
\end{enumerate}
\end{problem}

\begin{solution}
\begin{enumerate}
\item
方程 \(f(x)-x=0\) 为
\[
x^2+(a-1)x+a=0
\]
由于两根 \(x_1\)，\(x_2\) 都为正，韦达定理给出
\(x_1+x_2=1-a>0\)，\(x_1x_2=a>0\)
故 \(0<a<1\). 再由两根为不相等实根，
\[
\Delta=(a-1)^2-4a=a^2-6a+1>0
\]
结合 \(0<a<1\)，可得
\[
0<a<3-2\sqrt2
\]
反之，当 \(0<a<3-2\sqrt2\) 时，\(\Delta>0\)，且两根之和为正、积为正，所以两根均为正；又
\[
x_1+x_2=1-a<1
\]
从而两根均小于 \(1\). 因此
\[
0<a<3-2\sqrt2
\]
\item
又
\(f(0)=a\)，\(f(1)=1+2a\)
所以
\[
f(0)f(1)-f(0)=a(1+2a)-a=2a^2
\]
又因为 \(3-2\sqrt2<\frac{\sqrt2}{8}\)：两边均为正，且该不等式等价于 \(24<17\sqrt2\)，平方后为 \(576<578\). 因此
\[
a<\frac{\sqrt2}{8}
\]
从而
\[
2a^2<2\left(\frac{\sqrt2}{8}\right)^2=\frac1{16}
\]
故
\[
f(0)f(1)-f(0)<\frac1{16}
\]
\end{enumerate}
\end{solution}

\begin{problem}
已知数列 \(\{a_{n}\}\) 和 \(\{b_{n}\}\) 满足：\(a_{1} = 1\)，\(a_{2} = 2\)，\(a_{n} > 0\)，\(b_{n} = \sqrt{a_{n}a_{n + 1}}\)，\(\ n \in \mathbf{N}^{*}\)，且 \(\{b_{n}\}\) 是以 \(q\) 为公比的等比数列.
\begin{enumerate}
\item 证明：\( a_{n+2}=a_{n}q^{2} \).
\item 若 \(c_{n} = a_{2n - 1} + 2a_{2n}\)，证明数列 \(\{c_{n}\}\) 是等比数列.
\item 求和：\(\frac{1}{a_1} + \frac{1}{a_2} + \frac{1}{a_3} + \frac{1}{a_4} + \cdots + \frac{1}{a_{2n-1}} + \frac{1}{a_{2n}}\).
\end{enumerate}
\end{problem}

\begin{solution}
\begin{enumerate}
\item
由 \(a_n>0\) 及 \(b_n=\sqrt{a_na_{n+1}}\)，有
\[
b_n^2=a_na_{n+1}
\]
因为 \(\{b_n\}\) 是以 \(q\) 为公比的等比数列，且 \(b_n>0\)，所以 \(q>0\)，并且
\[
b_{n+1}=qb_n
\]
两边平方，得
\[
a_{n+1}a_{n+2}=q^2a_na_{n+1}
\]
因 \(a_{n+1}>0\)，可约去 \(a_{n+1}\)，从而
\[
a_{n+2}=a_nq^2
\]
\item
由此递推可得
\(a_{2n-1}=a_1q^{2n-2}=q^{2n-2}\)，\(a_{2n}=a_2q^{2n-2}=2q^{2n-2}\).
所以
\[
c_n=a_{2n-1}+2a_{2n}=5q^{2n-2}
\]
故 \(\{c_n\}\) 是首项为 \(5\)、公比为 \(q^2\) 的等比数列.
\item
最后，
\[
\sum_{k=1}^{2n}\frac1{a_k}
=\sum_{j=0}^{n-1}\left(\frac1{q^{2j}}+\frac1{2q^{2j}}\right)
=\frac32\sum_{j=0}^{n-1}q^{-2j}
\]
当 \(q^2\ne1\) 时，
\[
\sum_{k=1}^{2n}\frac1{a_k}
=\frac{3(q^{2n}-1)}{2q^{2n-2}(q^2-1)}
\]
当 \(q^2=1\) 时，
\[
\sum_{k=1}^{2n}\frac1{a_k}=\frac{3n}{2}
\]
\end{enumerate}
\end{solution}

\begin{problem}
在平面直角坐标系 \(xOy\) 中，过定点 \(C(0, p)\) 作直线与抛物线 \(x^{2} = 2py\)（\(p > 0\)）相交于 \(A\)，\(B\) 两点.

\begin{enumerate}
\item 若点 \(N\) 是点 \(C\) 关于坐标原点 \(O\) 的对称点，求 \(\triangle ANB\) 面积的最小值.
\item 是否存在垂直于 \(y\) 轴的直线 \(l\)，使得 \(l\) 被以 \(AC\) 为直径的圆截得的弦长恒为定值？若存在，求出 \(l\) 的方程；若不存在，说明理由.
\end{enumerate}

\FigureLayoutDeclare{img/2007/hubei_liberal/q21_fig1.png}{7.5cm}{151bp 220bp 123bp 46bp}
\bitmapfigure[max width=7.5cm]{img/2007/hubei_liberal/q21_fig1.png}
\end{problem}

\begin{solution}
\begin{enumerate}
\item
设过 \(C(0,p)\) 的直线方程为 \(y=kx+p\). 与抛物线 \(x^2=2py\) 联立，交点横坐标 \(x_1\)，\(x_2\) 满足
\[
x^2-2pkx-2p^2=0
\]
由韦达定理，
\(x_1+x_2=2pk\)，\(x_1x_2=-2p^2\).
又 \(N=(0,-p)\)，故
\[
S_{\triangle ANB}
=\frac12\left|x_1(kx_2+2p)-x_2(kx_1+2p)\right|
=p|x_1-x_2|
\]
由根的判别式，
\[
|x_1-x_2|
=\sqrt{(x_1+x_2)^2-4x_1x_2}
=\sqrt{4p^2k^2+8p^2}
=2p\sqrt{k^2+2}
\]
所以
\[
S_{\triangle ANB}=2p^2\sqrt{k^2+2}\geqslant2\sqrt2\,p^2
\]
当且仅当 \(k=0\) 时取等号. 因而三角形面积的最小值为 \(2\sqrt2\,p^2\).

\item
再设 \(l:y=t\). 以 \(AC\) 为直径的圆经过 \(A(x_1,kx_1+p)\) 和 \(C(0,p)\)，其方程为
\[
(x-x_1)x+(y-kx_1-p)(y-p)=0
\]
令 \(y=t\)，所得关于 \(x\) 的方程为
\[
x^2-x_1x+(t-kx_1-p)(t-p)=0
\]
其弦长恒定，当且仅当该方程的判别式与过 \(C\) 的直线斜率 \(k\) 无关. 因 \(A\) 在抛物线上，
\[
kx_1=\frac{x_1^2}{2p}-p
\]
所以判别式为
\[
\Delta=x_1^2-4(t-kx_1-p)(t-p)
=x_1^2\left(\frac{2t}{p}-1\right)-4t(t-p)
\]
点 \(A\) 可随过 \(C\) 的割线在抛物线上变化，故 \(x_1^2\) 不是常数. 要使 \(\Delta\) 恒定，必须且只需
\[
\frac{2t}{p}-1=0
\]
即 \(t=\frac p2\). 此时 \(\Delta=p^2\)，弦长为 \(\sqrt{\Delta}=p\)，确为定值. 因而存在所求直线
\[
l:y=\frac p2
\]
\end{enumerate}
\end{solution}
